% --- Encoding e font ---
\usepackage[utf8]{inputenc}
\usepackage[T1]{fontenc}
\usepackage{lmodern}

% --- Matematica ---
\usepackage{amsmath,amssymb,amsfonts,gensymb,amsthm}
\DeclareMathOperator*{\argmin}{arg\,min}
\DeclareMathOperator*{\argmax}{arg\,max}
\usepackage{centernot}
% --- Grafica, colori, float ---
\usepackage{graphicx}
\usepackage[dvipsnames]{xcolor}
\usepackage{float}
% --- Liste e ambienti ---
\usepackage{enumitem}
% --- Caption, numerazione figure ---
\usepackage{caption}
\usepackage{chngcntr}
% --- Caption, numerazione figure ---
\captionsetup{font=small, labelfont=bf, textfont=it}
% --- Indice/TOC personalizzazione (se ti serve) ---
\usepackage{tocloft}
% --- colorbox
\usepackage[most]{tcolorbox}
%--- spazio tra le riga
\usepackage{parskip}
\setlength{\parskip}{0.15cm}
% Elenchi compatti:
\setlist[itemize]{ topsep=2pt, itemsep=0pt, parsep=0pt}
\setlist[enumerate]{topsep=2pt, itemsep=0pt, parsep=0pt}

% Per il codice
\usepackage{listings}
% Definisci i colori esatti di VS Code (tema Light+)
\definecolor{vscodeblue}{RGB}{0,0,255}          % Keywords
\definecolor{vscodeorange}{RGB}{163,21,21}      % Strings
\definecolor{vscodegreen}{RGB}{0,128,0}         % Comments
\definecolor{vscodepurple}{RGB}{197,134,192}    % Control flow e import
\definecolor{vscodeyellow}{RGB}{121,94,38}      % Functions
\definecolor{vscodebg}{RGB}{255,255,255}        % Background
\definecolor{vscodeframe}{RGB}{200,200,200}     % Frame

% Configurazione Python stile VS Code
\lstdefinestyle{vscode}{
    language=Python,
    alsoletter={.},
    backgroundcolor=\color{vscodebg},
    commentstyle=\color{vscodegreen}\itshape,
    numberstyle=\tiny\color{gray},
    stringstyle=\color{vscodeorange},
    basicstyle=\ttfamily\small,
    breakatwhitespace=false,
    breaklines=true,
    captionpos=b,
    keepspaces=true,
    numbers=left,
    numbersep=8pt,
    showspaces=false,
    showstringspaces=false,
    showtabs=false,
    tabsize=4,
    frame=single,
    rulecolor=\color{vscodeframe},
    framesep=10pt,
    xleftmargin=15pt,
    framexleftmargin=15pt,
    % Resetta tutte le keyword e ridefiniscile
    morekeywords={},
    deletekeywords={for,if,while,else,elif,return,break,continue,try,except,finally,with,as,pass,raise,assert,import,from},
    % Keywords blu
    keywordstyle=\color{vscodeblue}\bfseries,
    morekeywords={def,class,lambda,and,or,not,is,in,del,global,nonlocal,yield,await,async,True,False,None,self},
    % Control flow VIOLA
    keywordstyle=[2]\color{vscodepurple}\bfseries,
    morekeywords=[2]{for,if,while,else,elif,return,break,continue,try,except,finally,with,as,pass,raise,assert,import,from},
    % Funzioni
    emph={sqrt,print,input,range,len,str,int,float,list,dict,set,tuple},
    emphstyle=\color{vscodeyellow}
}
\lstset{style=vscode}

%-- for definition (POI definisci defbox)
\newtcolorbox{defbox}[1]{
    colback=gray!10,
    colframe=gray!50,
    fonttitle=\bfseries,
    coltitle=black,
    title=\mbox{#1},  % Cambiato qui
    sharp corners,
    boxrule=0.5pt, breakable}
%-----
%-----
\usepackage{needspace}
%---- Notebox definition \begin{noteboxtitle}{title}
\newtcolorbox{noteboxtitle}[1]{
    colback=white,
    colframe=black!90,
    boxrule=0.8pt,
    arc=3mm,
    title=#1,
    fonttitle=\bfseries,
    left=5pt,
    right=5pt,
    top=5pt,
    bottom=5pt,
      breakable
}

% --- comment
\usepackage{comment}
% --- Margini ---
\usepackage[a4paper,
    top=2.5cm,
    bottom=2.5cm,
    left=2.5cm,
    right=2.5cm,
    bindingoffset=0.5cm
]{geometry}
% --- Titolo con immagine (opzionale) ---
\usepackage{titlepic} % solo se usi \titlepic
% --- Hyperref VA QUASI SEMPRE PER ULTIMO ---
\usepackage{hyperref}
\usepackage{adjustbox}
\usepackage{xltabular}
\usepackage{booktabs}
\raggedbottom


%---Per spezzare le tabelle
\usepackage{longtable}


% --- New structure for Example
\newtheoremstyle{smallexample}
  {3pt}% Space above
  {3pt}% Space below
  {\small}% Body font - testo dell'esempio piccolo
  {}% Indent amount
  {\small\bfseries}% Theorem head font - "Example 18.1" piccolo e grassetto
  {.}% Punctuation after theorem head
  {.5em}% Space after theorem head
  {}% Theorem head spec
\theoremstyle{smallexample}
\newtheorem{example}{Example}[section]
\usepackage[version=4]{mhchem}
\graphicspath{{Images/}}

 %% --- Definizione prima pagina 


%%___NEW COMMAND
\newcommand{\vs}{\vspace{0.2\baselineskip}}



%%%--- Librerie per grafici

\usepackage{tikz}
\usetikzlibrary{positioning, arrows.meta, shapes.geometric, calc}
\usepackage{tcolorbox}
\tcbuselibrary{breakable, skins}



\parindent 0px

